\documentclass[12pt, %
openright, 
oneside, %
%twoside, %TCC: Se seu texto tem mais de 100 páginas, descomente esta linha e comente a anterior
a4paper,    %
%english,   %
brazil]{facom-ufu-abntex2}


\usepackage{graphicx}
\graphicspath{{figuras/}{pictures/}{images/}{./}} % where to search for the images

\newcommand{\blue}[1]{\textcolor{blue}{#1}}
\newcommand{\red}[1]{\textcolor{red}{#1}}


\autor{José Lucas Ferreira de Lima} %TCC
\data{2025}
\orientador{Ronaldo Castro de Oliveira} %TCC
%\coorientador{Algum?} %TCC

% ---
% Informações de dados para CAPA e FOLHA DE ROSTO
% ---

\titulo{De voltas rápidas a decisões rápidas: o papel do big data na Fórmula 1} %TCC

\hypersetup{pdfkeywords={Big Data, Fórmula 1, telemetria, análise de dados, Kafka, UDP, Python, dashboard, machine learning, simulação, mensageria}}

\begin{document} 
\frenchspacing 

% ----------------------------------------------------------
% ELEMENTOS PRÉ-TEXTUAIS
% ----------------------------------------------------------
%\pretextual
\imprimircapa
\imprimirfolhaderosto


% ---
% Inserir folha de aprovação
% ---
%
% \includepdf{folhadeaprovacao_final.pdf} %TCC: depois de aprovado o trabalho, descomente esta linha e comente o próximo bloco para incluir scan da folha de aprovação.
%
\begin{folhadeaprovacao}

  \begin{center}
    {\ABNTEXchapterfont\large\imprimirautor}

    \vspace*{\fill}\vspace*{\fill}
    {\ABNTEXchapterfont\bfseries\Large\imprimirtitulo}
    \vspace*{\fill}
    
    \hspace{.45\textwidth}
    \begin{minipage}{.5\textwidth}
        \imprimirpreambulo
    \end{minipage}%
    \vspace*{\fill}
   \end{center}
    
   Trabalho aprovado. \imprimirlocal, 01 de novembro de 2016: %TCC:

   \assinatura{\textbf{\imprimirorientador} \\ Orientador}  
   \assinatura{\textbf{Professor}}% \\ Convidado 1} %TCC:
   \assinatura{\textbf{Professor}}% \\ Convidado 2} %TCC:
   %\assinatura{\textbf{Professor} \\ Convidado 3}
   %\assinatura{\textbf{Professor} \\ Convidado 4}
      
   \begin{center}
    \vspace*{0.5cm}
    {\large\imprimirlocal}
    \par
    {\large\imprimirdata}
    \vspace*{1cm}
  \end{center}
  
\end{folhadeaprovacao}
% ---


%%As seções dedicatória, agradecimento e epígrafe não são obrigatórias.
%%Só as mantenha se achar pertinente.

% ---
% Dedicatória
% ---
\begin{dedicatoria}
   \vspace*{\fill}
   \centering
   \noindent
   \textit{Dedico este trabalho aos meus pais, pelo apoio incondicional em toda a minha trajetória; aos meus amigos, 
   que foram pilares essenciais durante minha jornada acadêmica; e à minha avó, que nos deixou enquanto eu construía este trabalho, 
   mas cuja memória permanece viva em meu coração.}  %TCC:
   \vspace*{\fill}
\end{dedicatoria}
% ---

% ---
% Agradecimentos
% ---
%\begin{agradecimentos}
%Agradeço a \lipsum[30]. %TCC:
%\end{agradecimentos}
% ---

% ---
% Epígrafe
% ---
%\begin{epigrafe}
%    \vspace*{\fill}
%	\begin{flushright}
%		\textit{``Alguma citação que ache conveniente? \lipsum[10]''} %TCC:
%	\end{flushright}
%\end{epigrafe}
% ---



\begin{resumo} %TCC:

  A Fórmula 1 é um ambiente altamente tecnológico onde a análise de dados em 
  tempo real desempenha um papel crucial no desempenho das equipes, deixando a competitividade do esporte
  cada vez mais alta. Este trabalho explora a aplicação de \textbf{Big Data} na Fórmula 1, abordando tanto os conceitos teóricos 
  quanto uma implementação prática baseada na telemetria do jogo \textbf{F1 22}.
  
  O objetivo principal é demonstrar como a coleta e análise de dados podem fornecer insights valiosos para 
  a tomada de decisões estratégicas durante uma corrida. Para isso, foi desenvolvida uma aplicação em \textbf{Python} 
  que se conecta ao jogo pela rede local, estabelece um \textbf{socket} com este, recebe os dados de telemetria via \textbf{UDP}, 
  processa as informações, e envia os dados por mensageria utilizando \textbf{Kafka}. Os resultados são exibidos em um 
  \textbf{dashboard interativo} conectado a essa fila, permitindo análises em tempo real durante a corrida. 
  A aplicação foi testada em um ambiente controlado, e os resultados obtidos são discutidos em detalhes.
  
  A metodologia adotada inclui o estudo dos princípios do \textbf{Big Data} aplicados à F1, bem como o desenvolvimento 
  da pipeline de dados para captura e visualização dos dados do jogo. A implementação prática busca simular a experiência 
  real das equipes de F1 no uso de dados para análise de desempenho, ou seja, a visão do engenheiro de corrida durante a prova.

  Os resultados esperados incluem a demonstração do impacto do Big Data na performance das equipes e a 
  validação do potencial da tecnologia utilizada para fins de simulação e aprendizado. Este trabalho contribui 
  para a compreensão do uso da telemetria na F1 e reforça a importância da análise de dados em cenários de alta performance.

 \vspace{\onelineskip}
    
 \noindent
 \textbf{Palavras-chave}: big data, formula 1, engenharia de dados, python, pipeline. %TCC:
\end{resumo}

% ---
% inserir lista de ilustrações
% ---
\pdfbookmark[0]{\listfigurename}{lof}
\listoffigures*
\cleardoublepage
% ---

% ---
% inserir lista de tabelas
% ---
\pdfbookmark[0]{\listtablename}{lot}
\listoftables*
\cleardoublepage
% ---



% ---
% inserir lista de abreviaturas e siglas
% ---
\begin{siglas} %TCC:
  \item[Fig.] Area of the $i^{th}$ component
  \item[456] Isto é um número
  \item[123] Isto é outro número
  \item[Zézão] este é o meu nome
\end{siglas}
% ---

%% ---
%% inserir lista de símbolos, se for adequado ao trabalho. %TCC:
%% ---
%\begin{simbolos}
%  \item[$ \Gamma $] Letra grega Gama
%  \item[$ \Lambda $] Lambda
%  \item[$ \zeta $] Letra grega minúscula zeta
%  \item[$ \in $] Pertence
%\end{simbolos}
%% ---

% ---
% inserir o sumario
% ---
\pdfbookmark[0]{\contentsname}{toc}
\tableofcontents*
\cleardoublepage
% ---





% ----------------------------------------------------------
% ELEMENTOS TEXTUAIS
% ----------------------------------------------------------
\textual


% ----------------------------------------------------------
% Introdução
% ----------------------------------------------------------

\chapter[Introdução]{Introdução}
%TCC:
Contextualização, problema, hipótese, objetivo geral, objetivos específicos, justificativa e resultados esperados.
\section{Contextualização e Motivação}
\section{Problema de Pesquisa}
\section{Objetivos}
\subsection{Objetivo Geral}
\subsection{Objetivos Específicos}
\section{Justificativa}
\section{Metodologia do Trabalho}
\section{Estrutura do Trabalho}

% ---
% Revisão Bibliográfica
% ---

\chapter{Revisão Bibliográfica}
%TCC:
Um ou mais capítulos (por exemplo, se há duas linhas de trabalhos relacionados).
\section{Conceitos de Big Data}
\section{Aplicação de Big Data na Fórmula 1}
\section{Telemetria e Coleta de Dados em Tempo Real}
\section{Tecnologias Utilizadas}
\section{Trabalhos Relacionados}

% ---
% Desenvolvimento
% ---

\chapter{Desenvolvimento}
%TCC:
Um ou mais capítulos (por exemplo um para testes)
\section{Arquitetura da Solução}
\section{Coleta de Dados via Telemetria}
\section{Processamento e Armazenamento de Dados}
\section{Visualização dos Dados}
\section{Implementação e Testes}


\begin{figure}[!ht]
    \centering
	\includegraphics[width=0.55\linewidth]{imagemExemplo.pdf}
	\caption[Isso é o que aparece no sumário]{Imagem de exemplo.}
	\label{fig:graficosVariandoTamanhoRede}
\end{figure}

% ---
% Análise dos Resultados
% ---

\chapter{Análise dos Resultados}
\section{Validação da Coleta e Processamento}
\section{Desempenho e Eficiência da Aplicação}
\section{Limitações e Melhorias Futuras}




%TCC:
%TCC:
%TCC:
%TCC:

% ---
% Conclusão
% ---
\chapter{Conclusão e Trabalhos Futuros}
%TCC:
E daí?
\section{Revisão dos Objetivos e Contribuições}
\section{Desafios Encontrados}
\section{Aplicações Futuras}





% ----------------------------------------------------------
% ELEMENTOS PÓS-TEXTUAIS
% ----------------------------------------------------------
\postextual


% ----------------------------------------------------------
% Referências bibliográficas
% ----------------------------------------------------------
\bibliography{abntex2-modelo-references}


%% ----------------------------------------------------------
%% Apêndices TCC: só mantenha se for pertinente.
%% ----------------------------------------------------------

% ---
% Inicia os apêndices
% ---
\begin{apendicesenv}

% Imprime uma página indicando o início dos apêndices
\partapendices

% ----------------------------------------------------------
\chapter{Quisque libero justo}
% ----------------------------------------------------------

\lipsum[50]

% ----------------------------------------------------------
\chapter{Coisas que fiz e que achei interessante mas não tanto para entrar no corpo do texto}
% ----------------------------------------------------------
\lipsum[55-57]

\end{apendicesenv}
% ---


% ----------------------------------------------------------
% Anexos %TCC: so mantenha se pertinente.
% ----------------------------------------------------------

% ---
% Inicia os anexos
% ---
\begin{anexosenv}

% Imprime uma página indicando o início dos anexos
\partanexos

% ---
\chapter{Eu sempre quis aprender latim}
% ---
\lipsum[30]

% ---
\chapter{Coisas que eu não fiz mas que achei interessante o suficiente para colocar aqui}
% ---

\lipsum[31]

% ---
\chapter{Fusce facilisis lacinia dui}
% ---

\lipsum[32]

\end{anexosenv}

%---------------------------------------------------------------------
% INDICE REMISSIVO
%---------------------------------------------------------------------

\printindex



\end{document}